\chapter{Eta Invariants and Global Anomalies}
\label{ch:curved-eta-invariants-global-anomalies}

\ConventionsEuclidean

Chapter~\ref{ch:curved-background-gauge-fields-index-theory} treated the
local index density and the curvature of the determinant line.  This chapter
develops the complementary holonomy problem.  A local anomaly is detected by
infinitesimal motion in background-field space; a global anomaly is detected
by transporting the fermion functional integral around a closed loop of
backgrounds.  For free fermions the transport is computed by eta invariants
of Dirac operators on the associated mapping tori.

\section{Self-Adjoint Dirac Operators and Eta Functions}

Let \(Y\) be a closed oriented Riemannian spin manifold of odd dimension
\(2n+1\).  Let \(E\to Y\) be a Hermitian vector bundle with unitary connection
\(\nabla^E\).  The coupled Dirac operator
\[
  B_Y
  :
  C^\infty(Y,S_Y\otimes E)\longrightarrow
  C^\infty(Y,S_Y\otimes E)
\]
is self-adjoint on \(L^2(Y,S_Y\otimes E)\) when the Clifford action is chosen
skew-Hermitian in Euclidean signature.  Its spectrum is real, discrete, and
has finite-dimensional eigenspaces.  The eta function is initially defined by
\begin{equation}
  \eta_{B_Y}(s)
  =
  \sum_{\lambda\in\operatorname{Spec}(B_Y)\setminus\{0\}}
  \operatorname{sign}(\lambda)|\lambda|^{-s},
  \qquad
  \operatorname{Re}s\gg1 .
  \label{eq:eta-function-definition}
\end{equation}
For a Dirac-type operator the heat-kernel expansion gives a meromorphic
continuation to \(s\in\mathbb C\), regular at \(s=0\).  We define
\begin{equation}
  h(B_Y)=\dim\ker B_Y,
  \qquad
  \xi(B_Y)=\frac{\eta_{B_Y}(0)+h(B_Y)}2 .
  \label{eq:reduced-eta-invariant}
\end{equation}
The real number \(\xi(B_Y)\) depends on the metric and connection, but
\(\exp(2\pi\ii\xi(B_Y))\) is the natural phase variable: when eigenvalues
cross through zero, \(\xi\) changes by an integer.

\begin{remark}[Integer jump at a simple crossing]
\label{rem:eta-integer-jump}
Let \(B_t\), \(t\in(-\varepsilon,\varepsilon)\), be a smooth one-parameter
family of self-adjoint Dirac-type operators.  Suppose exactly one simple
eigenvalue \(\lambda(t)\) crosses zero, with
\(\lambda(t)<0\) for \(t<0\) and \(\lambda(t)>0\) for \(t>0\), and all other
eigenvalues stay away from zero.  Then the contribution of this eigenvalue to
\(\xi(B_t)\) jumps by \(+1\).  A crossing in the opposite direction gives a
jump by \(-1\).
Away from the crossing the eigenvalue contributes
\(\operatorname{sign}(\lambda)/2\) to
\((\eta(0)+h)/2\).  Just before the crossing this contribution is
\(-1/2\), at the crossing it is \(1/2\) because \(h=1\), and just after the
crossing it is \(1/2\).  Thus the one-sided jump from \(t<0\) to \(t>0\) is
\(+1\).  If the eigenvalue crosses from positive to negative, the
contribution changes from \(+1/2\) before the crossing to \(-1/2\) after the
crossing, again with \(h=1\) at the crossing instant, so the one-sided jump is
\(-1\).
\end{remark}

\begin{definition}[Spectral flow]
For a smooth path \(B_t\), \(0\leq t\leq1\), whose endpoints are invertible,
the spectral flow \(\operatorname{sf}(B_t)\) is the signed number of
eigenvalues crossing zero from negative to positive, counted with
multiplicity.  If the endpoints have kernels, one fixes a small perturbation
or uses the APS convention below; the resulting quantity is well-defined
modulo the endpoint-kernel bookkeeping appearing in \(\xi\).
\end{definition}

Remark~\ref{rem:eta-integer-jump} is the elementary reason that global
anomalies are naturally phases rather than real-valued functions of
\(\eta(0)\) alone.  The smooth part of \(\xi\) is local; the integer jumps
record spectral flow.

\section{The APS Boundary Formula}

Let \(X\) be a compact oriented Riemannian spin manifold of even dimension
\(2n+2\), with boundary \(\partial X=Y\).  Assume that, on a collar
\((-\epsilon,0]\times Y\), the metric and the connection have product form.
Let \(u\) be the inward normal coordinate.  In this collar the chiral Dirac
operator has the form
\begin{equation}
  D_X^+
  =
  c(\nu)\bigl(\partial_u+B_Y\bigr)
  \quad\hbox{up to the standard chiral identification,}
  \label{eq:aps-collar-dirac-form}
\end{equation}
where \(\nu\) is the outward unit normal and \(B_Y\) is the self-adjoint
tangential Dirac operator on \(Y\).

The APS boundary condition keeps the boundary values of spinors in the
negative spectral subspace of \(B_Y\):
\begin{equation}
  P_{\geq0}\bigl(\psi|_Y\bigr)=0,
  \label{eq:aps-boundary-condition}
\end{equation}
where \(P_{\geq0}\) is the orthogonal projection onto the span of eigenvectors
with eigenvalue \(\geq0\).  Including the zero eigenspace in \(P_{\geq0}\) is
the convention paired with the reduced invariant \(\xi=(\eta+h)/2\).

\begin{quotedtheorem}[Atiyah--Patodi--Singer index formula]
\label{qthm:aps-index-formula}
For the product-collar data above, the chiral Dirac operator with APS
boundary condition is Fredholm and
\begin{equation}
  \operatorname{Ind}(D_X^{+,\mathrm{APS}})
  =
  \int_X
  \bigl[\widehat A(TX)\operatorname{ch}(E)\bigr]_{2n+2}
  -
  \xi(B_Y).
  \label{eq:aps-index-formula}
\end{equation}
For non-product data one first deforms to product form near the boundary or
adds the corresponding local transgression term.  The sum of bulk density,
local boundary transgression, and eta term is the invariant integer.
\end{quotedtheorem}

The theorem is a theorem of global analysis.  In QFT it is used as the
mathematical statement that separates local anomaly density from spectral
boundary data.

\begin{proposition}[Cylinder variation of the reduced eta invariant]
\label{prop:eta-variation-local}
Let \(B_t\), \(0\leq t\leq1\), be a smooth path of odd-dimensional Dirac
operators obtained from a smooth path of metrics and unitary connections on a
fixed closed odd-dimensional spin manifold \(Y\).  Let
\[
  X=[0,1]\times Y
\]
carry the interpolating product-near-boundary data.  Then
\begin{equation}
  \xi(B_1)-\xi(B_0)
  \equiv
  \int_X
  \bigl[\widehat A(TX)\operatorname{ch}(E)\bigr]_{2n+2}
  \pmod{\mathbb Z}.
  \label{eq:eta-cylinder-variation}
\end{equation}
If the path has no zero crossings and the endpoint kernels have constant
dimension, differentiating this formula gives the usual local descent
variation of \(\xi\).
\end{proposition}

\begin{proof}
The boundary of \(X=[0,1]\times Y\) is
\[
  \partial X=Y_1\sqcup(-Y_0).
\]
The tangential operator on \(Y_1\) is \(B_1\), while the tangential operator on
the oppositely oriented component \(-Y_0\) is \(-B_0\).  Applying
\eqref{eq:aps-index-formula} to \(X\) gives
\[
  \operatorname{Ind}(D_X^{+,\mathrm{APS}})
  =
  \int_X
  \bigl[\widehat A(TX)\operatorname{ch}(E)\bigr]_{2n+2}
  -
  \xi(B_1)-\xi(-B_0).
\]
Since \(\eta_{-B_0}(0)=-\eta_{B_0}(0)\) and
\(h(-B_0)=h(B_0)\), one has
\[
  \xi(-B_0)
  =
  -\xi(B_0)+h(B_0).
\]
The Fredholm index and \(h(B_0)\) are integers.  Moving them to the other
side gives \eqref{eq:eta-cylinder-variation} modulo \(\mathbb Z\).
\end{proof}

Figure~\ref{fig:eta-cylinder-mapping-torus} records the cylinder/mapping-torus
geometry used in the eta-variation formula: interpolating between boundary
operators converts eta change into an APS index modulo integers.

\begin{figure}[t]
\centering
\begin{tikzpicture}[>=Stealth,scale=0.95]
  \draw[thick] (0,0) rectangle (4.2,2.2);
  \draw[->,thick] (0.45,1.1) -- (3.75,1.1);
  \node at (0.2,1.1) [left] {\(Y_0\)};
  \node at (4.4,1.1) [right] {\(Y_1\)};
  \node at (2.1,2.45) {\([0,1]\times Y\)};
  \node at (2.1,0.75) {\(B_t\)};
  \draw[rounded corners,thick] (6.0,0.15) rectangle (9.9,2.05);
  \draw[->,thick] (6.5,1.1) arc[start angle=180,end angle=-180,x radius=1.45,y radius=0.72];
  \node at (7.95,2.45) {mapping torus \(Y_\gamma\)};
  \node at (7.95,0.65) {\((x,1)\sim\gamma(x,0)\)};
\end{tikzpicture}
\caption{The APS cylinder computes the local change of \(\xi\) modulo
integers.  Identifying the two ends by a closed loop of backgrounds gives the
mapping torus whose eta invariant computes determinant-line holonomy.}
\label{fig:eta-cylinder-mapping-torus}
\end{figure}

\section{Determinant-Line Holonomy}

Let \(M\) be a closed even-dimensional spin manifold, and let
\(\mathcal B\) be a smooth parameter space of background metrics and
connections, before quotienting by large gauge transformations.  A family of
chiral Dirac operators
\[
  D_b^+:\Gamma(M,S^+\otimes E_b)\longrightarrow
  \Gamma(M,S^-\otimes E_b),
  \qquad b\in\mathcal B,
\]
defines the determinant line
\begin{equation}
  \operatorname{Det}D_b
  =
  \bigwedge\nolimits^{\mathrm{top}}\ker D_b^+
  \otimes
  \left(\bigwedge\nolimits^{\mathrm{top}}\ker D_b^-\right)^* .
  \label{eq:det-line-fiber}
\end{equation}
Where kernels jump this formula is not a literal vector bundle chart, but the
Quillen determinant construction glues these local descriptions to a smooth
complex line bundle over \(\mathcal B\).  The regulated chiral fermion
functional integral is a section of this line, not a function, until a
trivialization is chosen.

Let \(\gamma:S^1\to\mathcal B\) be a loop of backgrounds.  The mapping torus
\begin{equation}
  Y_\gamma
  =
  ([0,1]\times M)/\bigl((1,x)\sim(0,x)\hbox{ using }\gamma\bigr)
  \label{eq:mapping-torus-background-loop}
\end{equation}
is a closed odd-dimensional spin manifold with the induced bundle and
connection.  The Dirac operator on \(Y_\gamma\) will be denoted
\(B_{Y_\gamma}\).

\begin{quotedtheorem}[Bismut--Freed holonomy formula]
\label{qthm:bismut-freed-holonomy}
With the determinant-line and chirality conventions of
\eqref{eq:det-line-fiber}, the natural Bismut--Freed connection on
\(\operatorname{Det}D\) has holonomy
\begin{equation}
  \operatorname{Hol}_\gamma(\operatorname{Det}D)
  =
  \exp\!\left(-2\pi\ii\,\xi(B_{Y_\gamma})\right).
  \label{eq:bismut-freed-holonomy}
\end{equation}
Its curvature is the fiber integral over \(M\) of the degree
\((\dim M+2)\) component of
\(\widehat A(TM)\operatorname{ch}(E)\).
\end{quotedtheorem}

The curvature statement is the local anomaly polynomial.  The holonomy
statement is the global anomaly.  Therefore cancellation of the local anomaly
is the vanishing of the curvature of the total anomaly line, while
cancellation of the global anomaly is the triviality of the holonomy of that
flat line around every closed loop in background-field space.

\begin{proposition}[Finite gauge transformation as determinant-line transport]
\label{prop:finite-gauge-transformation-holonomy}
Let \(g\) be a gauge transformation of a background bundle on \(M\), and let
\(A_t\), \(0\leq t\leq1\), be a path from \(A\) to \(A^g\).  The ratio between
the transported determinant at \(A^g\) and the determinant at \(A\) is
computed by \eqref{eq:bismut-freed-holonomy} on the mapping torus
\begin{equation}
  Y_g
  =
  ([0,1]\times M)/\bigl((1,x)\sim(0,gx)\bigr).
  \label{eq:gauge-mapping-torus}
\end{equation}
If this phase is not \(1\), the chiral fermion measure does not descend to a
function on the quotient by that gauge transformation.
\end{proposition}

\begin{proof}
Parallel transport along \(A_t\) identifies
\(\operatorname{Det}D_A\) with \(\operatorname{Det}D_{A^g}\) by the
Bismut--Freed connection on the determinant line bundle over background
fields.  The gauge transformation \(g\) gives a canonical pullback
identification of \(D_{A^g}\) with \(D_A\), hence of their determinant lines.
Composing the parallel transport with this gauge identification gives an
automorphism of the one-dimensional line \(\operatorname{Det}D_A\), so it is
multiplication by a phase.  The corresponding closed path in the quotient
background groupoid is represented geometrically by gluing the two ends of
\([0,1]\times M\) with \(g\), namely by the mapping torus \(Y_g\).  The
Bismut--Freed holonomy formula applied to this closed family gives
\eqref{eq:bismut-freed-holonomy}.  A phase different from \(1\) means that
the determinant vector changes when the same quotient background is described
by the two representatives \(A\) and \(A^g\); hence the chiral fermion measure
does not descend to a function on the quotient.
\end{proof}

\section{Pfaffian Lines and Mod-Two Indices}

Some fermion systems have a real or pseudoreal structure that replaces the
complex determinant line by a Pfaffian line.  The four-dimensional
\(SU(2)\) fundamental Weyl fermion is the basic example.  The fundamental
representation is pseudoreal, and the spin representation supplies an
anti-linear symmetry that makes the Euclidean quadratic form skew on a real
vector space.  The path integral is then a Pfaffian.  Its possible global
anomaly is a sign.

\begin{definition}[Mod-two index]
\label{def:mod-two-index}
Let \(Y\) be a closed spin manifold of dimension \(8k+5\), and let \(E\) be a
pseudoreal bundle for which the coupled Dirac operator has the corresponding
real skew structure.  The mod-two index is
\begin{equation}
  \operatorname{Ind}_2(B_Y)
  =
  \dim\ker B_Y
  \pmod 2 .
  \label{eq:mod-two-index-definition}
\end{equation}
The parity is stable under continuous deformations preserving the real
structure.  The reason is that eigenvalues away from zero occur in pairs
\(\lambda,-\lambda\) for the real
skew structure relevant to the Pfaffian.  Under a deformation, nonzero
eigenvalues can enter or leave the kernel only through such paired crossings.
Thus the dimension of the kernel can change only by an even integer, and
\(\dim\ker B_Y\) modulo \(2\) is invariant.
\end{definition}

For such a Pfaffian line the Dai--Freed holonomy specializes to
\begin{equation}
  \operatorname{Hol}_\gamma(\operatorname{Pf}D)
  =
  (-1)^{\operatorname{Ind}_2(B_{Y_\gamma})}.
  \label{eq:pfaffian-mod-two-holonomy}
\end{equation}
This formula is the sign analogue of
\eqref{eq:bismut-freed-holonomy}.

\section{\texorpdfstring{The Witten \(SU(2)\) Global Anomaly}{The Witten SU(2) Global Anomaly}}

Let \(M=S^4\), and let \(g:S^4\to SU(2)\) represent the nontrivial element of
\[
  \pi_4(SU(2))\simeq\mathbb Z_2 .
\]
Form the mapping torus \(Y_g\) as in \eqref{eq:gauge-mapping-torus}.  A
single left-handed Weyl fermion in the fundamental representation of
\(SU(2)\) has Pfaffian holonomy
\begin{equation}
  \operatorname{Hol}_g(\operatorname{Pf}D)=-1 .
  \label{eq:witten-su2-fundamental-sign}
\end{equation}
Equivalently,
\begin{equation}
  \operatorname{Ind}_2(B_{Y_g,\square})=1 .
  \label{eq:witten-su2-fundamental-mod-two-index}
\end{equation}

\begin{quotedtheorem}[Mod-two index theorem for the Witten mapping torus]
\label{qthm:witten-su2-mod-two-index}
Let \(R_j\) be the irreducible \(SU(2)\) representation of isospin
\(j\in\frac12\mathbb Z_{\geq0}\), and set \(n=2j\).  With the trace-delta
normalization of Chapter~\ref{ch:curved-background-gauge-fields-index-theory},
\begin{equation}
  \operatorname{tr}_{R_j}(t^a t^b)
  =
  T_\Delta(R_j)\delta^{ab},
  \qquad
  T_\Delta(R_j)
  =
  \frac{n(n+1)(n+2)}6 .
  \label{eq:su2-trace-delta-dynkin-index}
\end{equation}
For the mapping torus of the generator of \(\pi_4(SU(2))\),
\begin{equation}
  \operatorname{Ind}_2(B_{Y_g,R_j})
  \equiv
  T_\Delta(R_j)
  \pmod2 .
  \label{eq:witten-su2-index-parity-dynkin}
\end{equation}
\end{quotedtheorem}

The theorem is a five-dimensional mod-two index theorem.  The algebraic
content of the anomaly criterion is elementary once
\eqref{eq:witten-su2-index-parity-dynkin} is known.

The isospin-\(j\) Weyl fermion has the original spin \(SU(2)\) global anomaly
precisely when \(2j\equiv1\pmod4\).  A collection of Weyl fermions in
representations \(R_{j_\alpha}\) is free of this anomaly precisely when
\begin{equation}
  \sum_\alpha T_\Delta(R_{j_\alpha})
  \equiv0
  \pmod2 .
  \label{eq:su2-global-anomaly-cancellation-condition}
\end{equation}
Write \(n=2j\).  The integer
\[
  T_\Delta(R_j)=\frac{n(n+1)(n+2)}6
\]
is the binomial coefficient \(\binom{n+2}{3}\).  Its parity is given by
Lucas' theorem: \(\binom{N}{3}\) is odd exactly when the binary digits of
\(3\) are contained in those of \(N\).  Since \(3\) has binary expansion
\(\cdots 0011\), this means \(N=n+2\equiv3\pmod4\), or
\[
  n\equiv1\pmod4 .
\]
Equation~\eqref{eq:witten-su2-index-parity-dynkin} then gives the Pfaffian
sign.  For several Weyl fermions the Pfaffian line is the tensor product of
the individual Pfaffian lines, so the signs multiply and the exponents add
modulo \(2\).

\begin{example}[First few \(SU(2)\) representations]
\label{ex:su2-global-anomaly-first-representations}
The sequence of trace-delta indices is
\[
\begin{array}{c|ccccc}
j & \frac12 & 1 & \frac32 & 2 & \frac52\\
\hline
T_\Delta(R_j) & 1 & 4 & 10 & 20 & 35 .
\end{array}
\]
Thus a single doublet has the anomaly, a single triplet does not, a
single four-dimensional representation does not, and a single
six-dimensional representation has the anomaly.  The vanishing of the
ordinary four-dimensional cubic \(SU(2)\) anomaly does not decide this
question because \(SU(2)\) has no independent symmetric cubic invariant.
\end{example}

\section{Dai--Freed Inflow}

The determinant or Pfaffian of a \(d\)-dimensional fermion is a vector in an
anomaly line over the background groupoid.  A \((d+1)\)-dimensional
invertible field theory trivializes this line when its boundary variation is
the inverse of the fermion-line variation.  In the free-fermion model the
closed \((d+1)\)-manifold partition function is the eta phase
\begin{equation}
  Z_{\mathrm{DF}}(Y)
  =
  \exp\!\left(2\pi\ii\,\xi(B_Y)\right)
  \label{eq:dai-freed-closed-phase}
\end{equation}
for a complex chiral determinant, with the corresponding mod-two sign in the
real Pfaffian case.

If \(Y\) has boundary \(M\), the APS construction assigns not a number but a
vector in the inverse determinant line over the boundary background on \(M\).
Pairing this vector with the boundary fermion determinant gives a number.
Changing the filling changes the trivialization by the closed-manifold phase
obtained by gluing the two fillings.  This is the precise geometric meaning
of anomaly inflow in the eta-invariant model.

\paragraph{Local descent as the infinitesimal limit of inflow.}
For a contractible loop of background fields, the Dai--Freed holonomy reduces
to the exponential of the descent integral determined by the local anomaly
polynomial.  For a noncontractible loop, the reduced eta invariant contains
additional flat anomaly-line holonomy data beyond the local polynomial.

If the loop is contractible, choose a disk of backgrounds filling it.  The
mapping torus bounds the total space over this disk.  Applying the APS index
formula to the bounding manifold expresses \(\xi\) modulo integers as the
integral of \(\widehat A\operatorname{ch}\).  The transgression of this
integral to the boundary loop is exactly the descent form.  If the loop is
not contractible in the background groupoid, such a disk need not exist; the
eta phase is then a genuine holonomy of a flat anomaly line.

\section{Interacting Theories and Anomaly Lines}

For free fermions, the anomaly line is constructed from a family of Dirac
operators.  For an interacting QFT, the same structure should be formulated
without choosing quadratic fields.  The background groupoid has objects given
by structured background fields and morphisms given by gauge transformations,
diffeomorphisms, and identifications of background data.  An anomaly line is
a functor from this groupoid to the groupoid of complex lines, equipped with
compatible gluing data.  The partition function of the anomalous theory is a
section of this functorial line.

Local counterterms change the trivialization of the line.  They shift the
connection one-form and the cocycle representative, but they do not change
the curvature class or the holonomy class.  Thus a genuine global anomaly is
not removed by changing local coordinates on the space of actions.

\begin{openproblem}[Global anomaly line for interacting QFT]
\label{op:global-anomaly-line-interacting-qft}
Construct the anomaly line of an interacting QFT over the space of background
fields, prove its local curvature and global holonomy formulae, and identify
the conditions under which it is represented by an invertible bulk field
theory.  The free-fermion determinant and Pfaffian lines developed above are
models for this structure, not a general construction of it.
\end{openproblem}
